% Source: https://artofproblemsolving.com/community/c3h568593_2009_national_target
% Stable identity: aops-topic-568593-q01
The two-digit primes less than $40$ are $11,13,17,19,23,29,31,37$, all without the digit $0$: there are $8$. The house number is determined by the ordered pair (first prime $AB$, second prime $CD$) with $AB\ne CD$: $8$ choices for $AB$ and then $7$ for $CD$, so there are $8\cdot7=\boxed{56}$ house numbers.

Multiplying by $2$ would count each house number twice. The product $8\cdot7$ already counts ordered pairs: for example $1317$ (with $AB=13$, $CD=17$) and $1713$ (with $AB=17$, $CD=13$) are both among the $56$. The factor $2$ would only be needed if we had first chosen an unordered pair of primes, $\dbinom82=28$ of them, and then decided their order: $28\cdot2=56$.
